\section{曲線の長さ}\label{節:曲線の弧長}

面積は長方形で,体積は柱体で近似した.では,曲線 $y=f(x)$ の長さはどう求めればよいだろうか.長方形は面積を測るための形なので,長さの近似には使えない.

すぐに思いつく候補を2つ挙げて,外しておこう.1つは横の幅 $\int_a^b\dd x=b-a$ だが,これは横方向の移動しか数えていない.もう1つは縦方向の変化の総量 $\int_a^b|f'(x)|\dd x$ だが,水平な線分では $f'=0$ となり,長さ $b-a$ を数えられない.横と縦の両方の移動を同時に数える必要がある.そこで,曲線上の近い2点を線分で結び,折れ線の長さで近似しよう.

実験として,\,$y=\frac23x^{\frac32}$ の $0\le x\le3$ の部分を考える.区間を $n$ 等分し,分点に対応する曲線上の点を順に線分で結んだ折れ線の長さを $L_n$ とする.三平方の定理で各線分の長さを求めて足すと,次の表のようになる.
\begin{center}
\begin{tabular}{ccccccc}
\toprule
$n$ & $1$ & $2$ & $4$ & $8$ & $16$ & $32$\\
\midrule
$L_n$ & $4.58258$ & $4.63181$ & $4.65396$ & $4.66245$ & $4.66535$ & $4.66627$\\
\bottomrule
\end{tabular}
\end{center}
分点を増やすと折れ線は曲線に近づき,\,$L_n$ は増えながら $4.667$ 付近に近づいていく.一方,横の幅は $3$,縦の変化の総量は $2\sqrt3\approx3.464$ であり,どちらも折れ線より短い.この極限の値を積分で表そう.

$f$ は $[a,b]$ で連続微分可能,すなわち導関数 $f'$ も連続であるとする.\,$[a,b]$ を $n$ 等分し,\,$\Delta x=\frac{b-a}{n}$,\,$x_k=a+k\Delta x$,\,$\Delta y_k=f(x_k)-f(x_{k-1})$ と置く.\,$k$ 番目の線分の長さは
\[
\sqrt{(\Delta x)^2+(\Delta y_k)^2}
 =\sqrt{1+\left(\frac{\Delta y_k}{\Delta x}\right)^2}\,\Delta x
\]
である.平均値の定理により,\,$x_{k-1}<\xi_k<x_k$ を満たすある $\xi_k$ で $\frac{\Delta y_k}{\Delta x}=f'(\xi_k)$ となるから,折れ線の長さは
\[
\sum_{k=1}^n\sqrt{1+f'(\xi_k)^2}\,\Delta x
\]
と書ける.これは連続関数 $\sqrt{1+f'(x)^2}$ のリーマン和なので,\,$n\to\infty$ として
\begin{equation}\label{式:グラフの弧長}
L=\int_a^b\sqrt{1+f'(x)^2}\dd x
\end{equation}
を得る.これを曲線の\textbf{弧長}という.根号の中の $1$ が横方向の移動を,\,$f'(x)^2$ が縦方向の変化を数えている.有限個の点で折れ曲がっていても,各部分が連続微分可能なら,部分ごとの長さを加えればよい.

実験の曲線に当てはめてみよう.\,$f(x)=\frac23x^{\frac32}$ より $f'(x)=\sqrt x$ であるから,
\begin{align*}
L&=\int_0^3\sqrt{1+x}\dd x\\
 &=\left[\frac23(1+x)^{\frac32}\right]_0^3\\
 &=\frac23(8-1)\\
 &=\frac{14}{3}
\end{align*}
となる.\,$\frac{14}{3}=4.666\cdots$ であり,表の値が近づいていく先と一致するのだ.

\begin{bookpoint}
\textbf{POINT}\par
$f'(x)$ を求めて $1+f'(x)^2$ を整理し,\,$\int_a^b\sqrt{1+f'(x)^2}\dd x$ を計算する.
\end{bookpoint}

\subsection{媒介変数表示}

円や,縦向きの接線を持つ曲線は,1つの関数 $y=f(x)$ のグラフとして扱いにくい.そこで
\[
x=x(t),\qquad y=y(t),\qquad \alpha\leq t\leq\beta
\]
と表された曲線の長さを考えよう.\,$x(t),y(t)$ は連続微分可能であるとする.

$t$ を $\Delta t$ だけ増やしたとき,横の変化はほぼ $x'(t)\Delta t$,縦の変化はほぼ $y'(t)\Delta t$ である.よって,対応する線分の長さはほぼ
\[
\sqrt{\{x'(t)\Delta t\}^2+\{y'(t)\Delta t\}^2}
 =\sqrt{x'(t)^2+y'(t)^2}\,\Delta t
\]
となる.これを足し合わせた折れ線の極限として
\begin{equation}\label{式:媒介曲線の弧長}
L=\int_\alpha^\beta\sqrt{x'(t)^2+y'(t)^2}\dd t
\end{equation}
を得る.導関数が連続であることから近似の誤差は区間全体で小さくできるが,ここではその評価の詳細は省き,この式を用いる.途中で $\sqrt{x'(t)^2+y'(t)^2}$ が $0$ になる点があっても,この式はそのまま使える.

ここで1つ断りをいれておこう.この積分は,\,$t$ に沿って点が進んだ長さを数える.同じ部分を2度たどれば2度数える.たとえば $x=\cos t$,\,$y=\sin t$ を $0\leq t\leq4\pi$ で積分すると $4\pi$ になるが,図形としての単位円の周の長さは $2\pi$ である.図形そのものの長さを求めるなら,1度だけたどる区間を選ぶ.

媒介変数表示は面積にも使える.\,$x(t)$ が単調増加し,\,$y(t)\geq0$ であるとする.このとき曲線は $x$ の関数のグラフであり,曲線と $x$ 軸の間の面積は $\int_{x(\alpha)}^{x(\beta)}y\dd x$ である.\,$x=x(t)$ と置換すると $\dd x=x'(t)\dd t$ なので,
\begin{equation}\label{式:媒介曲線の面積}
S=\int_\alpha^\beta y(t)x'(t)\dd t
\end{equation}
となる.\,$x(t)$ が途中で折り返すなら,区間を分けて,曲線が囲む部分を確認する.この式に絶対値を付けるだけで,いつでも図形の面積になるわけではない.縦の帯が重なれば,同じ面積を2度数えるからである.

\begin{bookpoint}
\textbf{POINT}\par
長さは $\sqrt{x'(t)^2+y'(t)^2}$ を,面積は $y(t)x'(t)$ を $t$ で積分する.根号を外すときは $t$ の範囲で符号を確かめる.
\end{bookpoint}

\begin{example}{媒介曲線の面積と長さ}
$a>0$ とする.曲線
\[
x=a(t-\sin t),\qquad y=a(1-\cos t),\qquad 0\leq t\leq2\pi
\]
の長さと,この曲線と $x$ 軸に囲まれた部分の面積を求めよ.
\end{example}

\paragraphbf{方針}
長さは式\eqref{式:媒介曲線の弧長}で求め,根号の中を半角の公式で2乗の形にする.面積は $x(t)$ が単調増加であることを確かめてから式\eqref{式:媒介曲線の面積}を使う.

\paragraphbf{解答}
長さ $8a$,面積 $3\pi a^2$

\paragraphbf{解法}
$x'(t)=a(1-\cos t)$,\,$y'(t)=a\sin t$ である.

\textbf{長さ}\quad 根号の中は
\begin{align*}
x'(t)^2+y'(t)^2
 &=a^2\{(1-\cos t)^2+\sin^2t\}\\
 &=a^2(2-2\cos t)\\
 &=4a^2\sin^2\frac t2
\end{align*}
となる.平方根は $2a\left|\sin\frac t2\right|$ である.\,$0\leq t\leq2\pi$ では $0\leq\frac t2\leq\pi$ なので $\sin\frac t2\geq0$ であり,
\begin{align*}
L&=2a\int_0^{2\pi}\sin\frac t2\dd t\\
 &=\left[-4a\cos\frac t2\right]_0^{2\pi}\\
 &=4a-(-4a)\\
 &=8a.
\end{align*}

\textbf{面積}\quad $x'(t)=a(1-\cos t)$ は $0<t<2\pi$ で正なので,\,$x$ は区間全体で単調増加する.また $y\geq0$ で,両端では $y=0$ である.よって
\begin{align*}
S&=\int_0^{2\pi}a(1-\cos t)\cdot a(1-\cos t)\dd t\\
 &=a^2\int_0^{2\pi}(1-2\cos t+\cos^2t)\dd t\\
 &=a^2\int_0^{2\pi}\left(\frac32-2\cos t+\frac12\cos2t\right)\dd t\\
 &=a^2\cdot\frac32\cdot2\pi\\
 &=3\pi a^2.
\end{align*}
$\cos t$ と $\cos2t$ の項は,\,$0$ から $2\pi$ まで積分すると $0$ になる.
%
